% generado por gen_tips.py; no editar
\makeatletter
\expandafter\gdef\csname CGTip@sec:intro\endcsname{Section 1: Introduction (p. 2)}
\expandafter\gdef\csname CGTip@subsec:clock_intro\endcsname{Section 1.1: The clock (p. 2)}
\expandafter\gdef\csname CGTip@eq:intro_RvM\endcsname{Eq. (1.1), p. 2: N(t)=Nz(t)+S(t)+(varpi(t))/pi , Nz(t)=t/2pi logt/(2pi e) +7/8 , S(t)=1/pi argzeta (1/2 +it ),}
\expandafter\gdef\csname CGTip@eq:intro_clock\endcsname{Eq. (1.2), p. 2: Nz(tau\_k)=k-1/2 .}
\expandafter\gdef\csname CGTip@subsec:results\endcsname{Section 1.2: Results (p. 3)}
\expandafter\gdef\csname CGTip@eq:admissible\endcsname{Eq. (1.3), p. 3: zeta(sigma+it)<<\_varepsilon|t|\textasciicircum{}mathfrak e(sigma)+varepsilon uniformly for 1/2 <=sigma<=1, |t|>=2 .}
\expandafter\gdef\csname CGTip@eq:reflected\endcsname{Eq. (1.4), p. 3: 1/2 -sigma, sigma<0, 1/2 -sigma+mathfrak e(1-sigma), 0<=sigma<1/2 , mathfrak e(sigma), 1/2 <=sigma<=1, 0, sigma>1,endcases}
\expandafter\gdef\csname CGTip@thm:envelope\endcsname{Theorem 1.1 (Envelope) (p. 3): For every admissible profile mathfrak e, beginequation mathfrak a<=Lambda(mathfrak e):=1/2pilog2 int\_0\textasciicircum{}1(mathfrak e\textasciicircum{}*(sigma))/(sqrt(sigma(2-sigma))) dsigma . endequation}
\expandafter\gdef\csname CGTip@eq:envelope_formula\endcsname{Eq. (1.5), p. 3: mathfrak a<=Lambda(mathfrak e):=1/2pilog2 int\_0\textasciicircum{}1(mathfrak e\textasciicircum{}*(sigma))/(sqrt(sigma(2-sigma))) dsigma .}
\expandafter\gdef\csname CGTip@eq:bourgain_nodes\endcsname{Eq. (1.6), p. 3: (1/2 ,13/84 ), (139/194 ,13/194 ), (173/207 ,13/414 ), (773/854 ,13/854 ), (820/867 ,13/1734 ), (3387/3494 ,13/3494 ), (1149/1169 ,13/7014 ), (13921/14054 ,13/14054 ), (1,0). endgathered}
\expandafter\gdef\csname CGTip@cor:bourgain\endcsname{Corollary 1.2 (Bourgain profile) (p. 3): mathfrak a<=Lambda(eB)=0.0968160063ldots}
\expandafter\gdef\csname CGTip@thm:explicit\endcsname{Theorem 1.3 (Explicit bound) (p. 3): For every T>= e, beginalign |S(T)| <=0.09783log T+0.7214loglog T+10.6, |N(T)-Nz(T) | <=0.09783log T+0.7214loglog T+10.7. endalign}
\expandafter\gdef\csname CGTip@eq:explicit_S\endcsname{Eq. (1.7), p. 3: |S(T)| <=0.09783log T+0.7214loglog T+10.6, |N(T)-Nz(T) | <=0.09783log T+0.7214loglog T+10.7.}
\expandafter\gdef\csname CGTip@eq:explicit_N\endcsname{Eq. (1.8), p. 3: |S(T)| <=0.09783log T+0.7214loglog T+10.6, |N(T)-Nz(T) | <=0.09783log T+0.7214loglog T+10.7.}
\expandafter\gdef\csname CGTip@cor:clock\endcsname{Corollary 1.4 (The clock) (p. 3): beginenumerate[label=(alph*)] item limsup\_k->inf|tau\_k-gamma\_k|<=2piLambda(eB)=0.6083129ldots<0.60832 . item For every k>=1, with mathfrak m\_k:=min(tau\_k,gamma\_k)>2pi, beginequation |tau\_k-gamma\_k|<=(2pi (0.09783loggamma\_k+0.7214logloggamma\_k+10.2 ))/(log(mathfrak m\_k/2pi)) . endequation endenumerate}
\expandafter\gdef\csname CGTip@eq:clock_explicit\endcsname{Eq. (1.9), p. 3: |tau\_k-gamma\_k|<=(2pi (0.09783loggamma\_k+0.7214logloggamma\_k+10.2 ))/(log(mathfrak m\_k/2pi)) .}
\expandafter\gdef\csname CGTip@eq:Euler_phase\endcsname{Eq. (1.10), p. 4: Ph\_m(t)=-1/pi sum\_nu=1\textasciicircum{}marg (1-p\_nu\textasciicircum{}-1/2-it )=-1/pi sum\_nu=1\textasciicircum{}msum\_r>=1(sin(rtlog p\_nu))/(r p\_nu\textasciicircum{}r/2)}
\expandafter\gdef\csname CGTip@eq:F_def\endcsname{Eq. (1.11), p. 4: Ph(t)= (2-t/(2\textasciicircum{}j) )Ph\_m\_j(t)+ (t/(2\textasciicircum{}j) -1 )Ph\_m\_j+1(t), Fc(t)=Nz(t)+Ph(t).}
\expandafter\gdef\csname CGTip@eq:xk_def\endcsname{Eq. (1.12), p. 4: x\_k=int\_0\textasciicircum{}infInd\_Fc(t)<k-1/2 dt (k>=1).}
\expandafter\gdef\csname CGTip@thm:sublevel\endcsname{Theorem 1.5 (Localisation by primes) (p. 4): The numbers x\_k are finite and 0<x\_1<x\_2<*s. For k>=2 (equivalently tau\_k>=2pi e, since Nz(2pi e)=7/8 ), x\_k is determined by k and by the primes p\_nu with nu<=lfloor2sqrt(tau\_k)rfloor; no ordinate and no value of zeta enters its definition. There are absolute constants C and T\_0 such that, with Kc\_T=k:T<gamma\_k<=2T, ...}
\expandafter\gdef\csname CGTip@eq:main_sublevel\endcsname{Eq. (1.13), p. 4: sum\_kinKc\_T(x\_k-gamma\_k)\textasciicircum{}2<= C T/(log T) , hence 1/(no.Kc\_T) sum\_kinKc\_T(x\_k-gamma\_k)\textasciicircum{}2<=(C')/(log\textasciicircum{}2T) (T>= T\_0),}
\expandafter\gdef\csname CGTip@cor:proportion\endcsname{Corollary 1.6 (p. 4): For every mathsf A>0 and T>= T\_0, no.kinKc\_T:|x\_k-gamma\_k|>mathsf A*2pi/(log T) <=(C'')/(mathsf A\textasciicircum{}2) no.Kc\_T, with C'' absolute.}
\expandafter\gdef\csname CGTip@subsec:literature\endcsname{Section 1.3: Relation to the literature (p. 4)}
\expandafter\gdef\csname CGTip@sec:notation\endcsname{Section 2: Notation (p. 7)}
\expandafter\gdef\csname CGTip@sec:external\endcsname{Section 3: Results from the literature (p. 9)}
\expandafter\gdef\csname CGTip@thm:known_zeta\endcsname{Theorem 3.3 (Explicit bounds for zeta) (p. 9): For real t: beginenumerate[label=(alph*)] item |zeta(1/2 +it)|<=0.618|t|\textasciicircum{}1/6log|t| for 3<|t|<= e\textasciicircum{}105 [HiaryPatelYang2024], and |zeta(1/2 +it)|<=66.7|t|\textasciicircum{}27/164 for |t|>e\textasciicircum{}105 [PatelYang2023]; item |zeta(sigma\_m+it)|<=1.546|t|\textasciicircum{}1/(2\textasciicircum{}m-2)log|t| for |t|>=3 and every integer m>=4, sigma\_m=1-m/(2\textasciicircum{}m-2) [Yang2024] (see Remark ...}
\expandafter\gdef\csname CGTip@eq:backlund_f\endcsname{Eq. (3.1), p. 9: mathfrak f\_n(s)=1/2 [ ((s+iT-1)zeta(s+iT) )\textasciicircum{}n+ ((s-iT-1)zeta(s-iT) )\textasciicircum{}n ],}
\expandafter\gdef\csname CGTip@eq:backlund_Z\endcsname{Eq. (3.2), p. 9: mathcal Z\_n(c,R)=1/n sum\_wlogR/(|w-c|) ,}
\expandafter\gdef\csname CGTip@thm:known_exponents\endcsname{Theorem 3.7 (Exponent pairs) (p. 10): beginenumerate[label=(alph*)] item If (kappa,lambda) is an exponent pair, so is A(kappa,lambda)= (kappa/(2kappa+2) ,1/2 +lambda/(2kappa+2) ) [Ivic1985]. item If (kappa,lambda) is an exponent pair and 1/2 <=sigma<=lambda-kappa, then zeta(sigma+it)<<|t|\textasciicircum{}(kappa+lambda-sigma)/2log|t| for |t|>=2 [Ivic1985]; in particular ...}
\expandafter\gdef\csname CGTip@sec:explicit_lemmas\endcsname{Section 4: Explicit lemmas (p. 11)}
\expandafter\gdef\csname CGTip@lem:varpi\endcsname{Lemma 4.1 (p. 11): For t>0, vartheta(t)=t/2 logt/(2pi e) -pi/8 +varpi(t) with |varpi(t)|<=0.53/t. Consequently [eq:intro\_RvM] holds for t>0 not an ordinate.}
\expandafter\gdef\csname CGTip@lem:chi\endcsname{Lemma 4.3 (p. 11): For -0.01<=sigma<=1/2 and t>=98, |chi(sigma+it)|<=1.0070 (t/2pi)\textasciicircum{}1/2-sigma.}
\expandafter\gdef\csname CGTip@eq:stirling_modulus\endcsname{Eq. (4.2), p. 12: log|Gamma(z)|=(omega-1/2 )log y-pi/2 y+1/2 log2pi+phi(omega,y)+Remathcal R(z), phi(omega,y)=(omega-1/2 )/2 log (1+(omega\textasciicircum{}2)/(y\textasciicircum{}2) )+yarctanomega/y -omega . endgathered}
\expandafter\gdef\csname CGTip@subsec:explicit_profile\endcsname{Section 4.1: The explicit profile (p. 12)}
\expandafter\gdef\csname CGTip@eq:explicit_nodes\endcsname{Eq. (4.3), p. 12: (1/2 ,27/164 ), (sigma\_m,1/(2\textasciicircum{}m-2) ) (4<= m<=9), (1,0),}
\expandafter\gdef\csname CGTip@lem:explicit_zeta\endcsname{Lemma 4.4 (p. 12): For -0.01<=sigma<=3 and |t|>=98, |zeta(sigma+it)|<=2.74 |t|\textasciicircum{}eE\textasciicircum{}*(sigma)log|t|.}
\expandafter\gdef\csname CGTip@eq:boundary\endcsname{Eq. (4.4), p. 12: |Xi(varsigma\_i+it)|<=2|Q+s|\textasciicircum{}1+e\_ilog|Q+s| (s=varsigma\_i+it).}
\expandafter\gdef\csname CGTip@eq:right_profile\endcsname{Eq. (4.5), p. 12: |zeta(s)| <=2 (1+84.5/(98\textasciicircum{}2) )\textasciicircum{}1+27/164 (1+(84.5/98\textasciicircum{}2+pi/2)/log98 )t\textasciicircum{}eE\textasciicircum{}*(sigma)log t <=2.7167 t\textasciicircum{}eE\textasciicircum{}*(sigma)log t (1/2 <=sigma<=3). endsplit}
\expandafter\gdef\csname CGTip@rem:robust\endcsname{Remark 4.5 (Robustness of the inputs) (p. 13): Theorem [thm:known\_zeta] enters Theorem [thm:explicit] only through [eq:boundary], that is, through bounds with constant 2 on the lines sigma=varsigma\_i, and each input has room. (i) On sigma=1/2 , the bound of [HiaryPatelYang2024] alone gives [eq:boundary] for 3<|t|<= e\textasciicircum{}577, since |t|\textasciicircum{}1/492<= e\textasciicircum{}577/492<2/0.618; above ...}
\expandafter\gdef\csname CGTip@sec:backlund\endcsname{Section 5: Backlund's method on a disc near 1 (p. 13)}
\expandafter\gdef\csname CGTip@prop:master\endcsname{Proposition 5.1 (p. 13): Let c,R satisfy [eq:backlund\_conditions] with c<=2 and R<=2, and let T>=3 be a non-ordinate. Suppose that, for a function mathfrak h>=0 on [c-R,c+R], a constant K>=1 and a function Omega>=1 on [T-2,T+2] such that tmapsto t\textasciicircum{}mathfrak h(sigma)Omega(t) is nondecreasing there for every sigma, beginequation ...}
\expandafter\gdef\csname CGTip@eq:zeta_hyp\endcsname{Eq. (5.1), p. 13: |zeta(sigma+it)|<= K t\textasciicircum{}mathfrak h(sigma)Omega(t) (c-R<=sigma<= c+R, T-2<= t<= T+2).}
\expandafter\gdef\csname CGTip@eq:master\endcsname{Eq. (5.2), p. 13: |S(T)|<= 1/(2mathfrak d) [log(zeta(c))/(zeta(2c)) +log(T+3)/T +(mathcal I\_mathfrak h)/pi log(T+2)+logOmega(T+2)+log K ] +(logzeta(sigma\textasciicircum{}sharp))/pi +(sigma\textasciicircum{}sharp-1/2 )/(pi T) +(mathcal E(T,delta\textasciicircum{}sharp))/2pi +1/4 , endsplit}
\expandafter\gdef\csname CGTip@rem:lindelof\endcsname{Remark 5.2 (p. 14): Admissible profiles are nonnegative by definition, so mathfrak e\textasciicircum{}*>=0\textasciicircum{}* pointwise, where 0\textasciicircum{}*(sigma)=max(1/2 -sigma,0) is the reflected extension of the zero profile, and [ Lambda(mathfrak e)>=Lambda(0)=1/2pilog2 int\_0\textasciicircum{}1/2((1/2-sigma) dsigma)/(sqrt(sigma(2-sigma))) =(3sqrt3-pi)/12pilog2 =0.0786252ldots ] The Lindel"of ...}
\expandafter\gdef\csname CGTip@sec:clock\endcsname{Section 6: The clock (p. 16)}
\expandafter\gdef\csname CGTip@lem:s_bounds\endcsname{Lemma 6.1 (p. 16): For every k>=1, G\_k\textasciicircum{}-+1/2 <= s\_k<= G\_k\textasciicircum{}+-1/2 ; in particular |s\_k|<=max(|G\_k\textasciicircum{}-|,|G\_k\textasciicircum{}+|)-1/2 .}
\expandafter\gdef\csname CGTip@prop:clock_sharp\endcsname{Proposition 6.2 (The envelope of the clock) (p. 16): limsup\_k->inf|tau\_k-gamma\_k|=2pimathfrak a. In particular Corollary [cor:clock](a) is equivalent to Corollary [cor:bourgain], and the Lindel"of hypothesis would give limsup\_k|tau\_k-gamma\_k|<=2piLambda(0)=0.49401ldots (Remark [rem:lindelof]).}
\expandafter\gdef\csname CGTip@sec:sublevel\endcsname{Section 7: Localisation by primes (p. 17)}
\expandafter\gdef\csname CGTip@subsec:construction\endcsname{Section 7.1: The construction (p. 17)}
\expandafter\gdef\csname CGTip@eq:phasebound\endcsname{Eq. (7.1), p. 17: |Ph\_m(t)|<=m/4 .}
\expandafter\gdef\csname CGTip@eq:globalphasebound\endcsname{Eq. (7.2), p. 17: |Ph(t)|<=1/4 sqrt(2t) (t>0).}
\expandafter\gdef\csname CGTip@prop:ordered\endcsname{Proposition 7.1 (p. 17): For every k>=1, x\_k is finite and 0<x\_k<x\_k+1. If tau\_k>=2pi e, the integral [eq:xk\_def] may be truncated at 2tau\_k, and its integrand uses only the primes p\_nu with nu<=lfloor2sqrt(tau\_k)rfloor.}
\expandafter\gdef\csname CGTip@subsec:transfer\endcsname{Section 7.2: A deterministic transfer inequality (p. 17)}
\expandafter\gdef\csname CGTip@lem:maximal\endcsname{Lemma 7.2 (p. 17): Let A<= B be integers and D=(D\_ell)\_A<=ell<= B nonnegative. For A<=ell<= B let (M D)\_ell=max1/(b-a+1) sum\_r=a\textasciicircum{}bD\_r: A<= a<=ell<= b<= B. Then sum\_ell=A\textasciicircum{}B(M D)\_ell\textasciicircum{}3<=24sum\_ell=A\textasciicircum{}BD\_ell\textasciicircum{}3.}
\expandafter\gdef\csname CGTip@thm:transfer\endcsname{Theorem 7.3 (Sublevel transfer) (p. 18): In this theorem (gamma\_k), N and D\_ell are general: let (gamma\_k)\_k>=1 be any nondecreasing real sequence with finitely many terms in each bounded set, N(t)=no.k:gamma\_k<= t, and a<b with ell\_0=N(a)>=1, ell\_1=N(b). Put D\_ell=gamma\_ell+1-gamma\_ell for ell\_0<=ell<=ell\_1. Let f be a measurable real function on [a,b], and ...}
\expandafter\gdef\csname CGTip@eq:transfer\endcsname{Eq. (7.3), p. 18: sum\_a<gamma\_k<= b(x\textasciicircum{}f\_k-gamma\_k)\textasciicircum{}2<=2*24\textasciicircum{}1/3 (sum\_ell=ell\_0\textasciicircum{}ell\_1D\_ell\textasciicircum{}3 )\textasciicircum{}1/2 (int\_a\textasciicircum{}bq\_f\textasciicircum{}4 dt )\textasciicircum{}1/2,}
\expandafter\gdef\csname CGTip@eq:bathtub\endcsname{Eq. (7.4), p. 18: (x\textasciicircum{}f\_k-gamma\_k)\textasciicircum{}2<= u\_k\textasciicircum{}2+v\_k\textasciicircum{}2<=2int\_E\_k|t-gamma\_k| dt .}
\expandafter\gdef\csname CGTip@rem:quantile\endcsname{Remark 7.4 (Quantile form) (p. 19): In this remark mathsf m\_f and mathsf m\_N are local. Let mathsf m\_f be the image of Lebesgue measure on [a,b] under f, a measure of mass b-a on R, and mathsf m\_N the image under N, which is the sum of the point masses meastin[a,b]:N(t)=ell delta\_ell, ell\_0<=ell<=ell\_1. Then x\textasciicircum{}f\_k-a=mathsf m\_f ((-inf,k-1/2 ) ) and ...}
\expandafter\gdef\csname CGTip@subsec:arith_inputs\endcsname{Section 7.3: The arithmetic inputs (p. 19)}
\expandafter\gdef\csname CGTip@lem:gaps\endcsname{Lemma 7.5 (Third moment of the gaps) (p. 19): sum\_T<gamma\_ell<=2T(gamma\_ell+1-gamma\_ell)\textasciicircum{}3<< T/log\textasciicircum{}2T for T>= T\_0.}
\expandafter\gdef\csname CGTip@lem:fourth\endcsname{Lemma 7.6 (Fourth moment of the counting error) (p. 19): For U=2\textasciicircum{}j with j>= j\_0, int\_U\textasciicircum{}2U|Fc(t)-N(t)|\textasciicircum{}4dt<< U.}
\expandafter\gdef\csname CGTip@eq:phase_split\endcsname{Eq. (7.5), p. 20: Ph\_m(t)=1/pi Imsum\_p<Yp\textasciicircum{}-1/2-it+1/pi Immathscr D\_1(t)+1/2pi Immathscr D\_2(t)+Ph\textasciicircum{}(3)(t), mathscr D\_1(t)=sum\_Y<= p<= Xp\textasciicircum{}-1/2-it, mathscr D\_2(t)=sum\_p<= Xp\textasciicircum{}-1-2it, endgathered}
\expandafter\gdef\csname CGTip@rem:bare_clock\endcsname{Remark 7.7 (p. 21): The bare clock tau\_k is the case Ph=0. Then q\_f is governed by S, Selberg's moments give int\_T\textasciicircum{}2TS\textasciicircum{}4asymp T(loglog T)\textasciicircum{}2 [Selberg1946a] (see [Titchmarsh1986]), and [eq:transfer] (with the window [T/2,3T] localised exactly as in the proof of Theorem [thm:sublevel], now with Ph=0) only yields ...}
\expandafter\gdef\csname CGTip@sec:conclusion\endcsname{Section 8: Conclusion (p. 21)}
\expandafter\gdef\csname CGTip@app:code\endcsname{Section A: Computation of the constants (p. 22)}
\expandafter\gdef\csname CGTip@app:lean\endcsname{Section B: Formal verification in Lean 4 (p. 27)}
\expandafter\gdef\csname CGTip@app:sources\endcsname{Section C: Correspondence with the sources (p. 28)}
\expandafter\gdef\csname CGBib@AlpogeFurmanZeta23\endcsname{L. Alp"oge, R. Furman, More than two thirds of the zeta zeros are simple and on the critical line, arXiv:2608.13637v2, 2026, doi:10.48550/arXiv.2608.13637 (the argument is attributed there to Claude, Anthropic); Lean 4 formalisation Zeta23, Apache-2.0, https://github.com/anthropics/zeta-23-lean.}
\expandafter\gdef\csname CGBib@Backlund1918\endcsname{R.J. Backlund, "Uber die Nullstellen der Riemannschen Zetafunktion, Acta Math. 41 (1916) 345--375, doi:10.1007/BF02422950.}
\expandafter\gdef\csname CGBib@Bellotti2024\endcsname{C. Bellotti, Explicit bounds for the Riemann zeta function and a new zero-free region, J. Math. Anal. Appl. 536 (2024) 128249, doi:10.1016/j.jmaa.2024.128249.}
\expandafter\gdef\csname CGBib@BellottiWong2025\endcsname{C. Bellotti, P.-J. Wong, Improved estimates for the argument and zero-counting function of the Riemann zeta-function, Math. Comp. (2025, published online 14 November 2025), statements numbered as in arXiv:2412.15470v2, doi:10.1090/mcom/4133.}
\expandafter\gdef\csname CGBib@BennettEtAl2021\endcsname{M.A. Bennett, G. Martin, K. O'Bryant, A. Rechnitzer, Counting zeros of Dirichlet L-functions, arXiv:2005.02989v1, 2020, cited with its numbering, doi:10.48550/arXiv.2005.02989; published version: Math. Comp. 90 (2021) 1455--1482, doi:10.1090/mcom/3599.}
\expandafter\gdef\csname CGBib@Bourgain2017\endcsname{J. Bourgain, Decoupling, exponential sums and the Riemann zeta function, J. Amer. Math. Soc. 30 (2017) 205--224, theorems numbered as in arXiv:1408.5794v2, doi:10.1090/jams/860.}
\expandafter\gdef\csname CGBib@CasanaPaperMath2026\endcsname{V. Casa na-Giner, PaperMath: Lean 4 formalisations accompanying the author's articles, Git repository, https://github.com/vcasana1/PaperMath, commit 97f4e7c, 2026.}
\expandafter\gdef\csname CGBib@DLMF\endcsname{NIST Digital Library of Mathematical Functions, S5.11(ii), Release 1.2.4 of 2025-03-15, F.W.J. Olver et al., eds., https://dlmf.nist.gov/5.11.ii.}
\expandafter\gdef\csname CGBib@FrancaLeClair2015\endcsname{G. Francca, A. LeClair, Transcendental equations satisfied by the individual zeros of Riemann zeta, Dirichlet and modular L-functions, Commun. Number Theory Phys. 9 (2015) 1--50, doi:10.4310/CNTP.2015.v9.n1.a1.}
\expandafter\gdef\csname CGBib@Fiori2025\endcsname{A. Fiori, A note on the Phragm'en--Lindel"of theorem, arXiv:2502.13282v3, 2026, doi:10.48550/arXiv.2502.13282.}
\expandafter\gdef\csname CGBib@Fujii1974\endcsname{A. Fujii, On the zeros of Dirichlet L-functions. I, Trans. Amer. Math. Soc. 196 (1974) 225--235, doi:10.1090/S0002-9947-1974-0349603-2.}
\expandafter\gdef\csname CGBib@Fujii1981\endcsname{A. Fujii, On the zeros of Dirichlet L-functions. II (with corrections to ``On the zeros of Dirichlet L-functions. I'' and the subsequent papers), Trans. Amer. Math. Soc. 267 (1981) 33--40, doi:10.1090/S0002-9947-1981-0621970-5.}
\expandafter\gdef\csname CGBib@Gonek2012\endcsname{S.M. Gonek, Finite Euler products and the Riemann hypothesis, Trans. Amer. Math. Soc. 364 (2012) 2157--2191, doi:10.1090/S0002-9947-2011-05546-7.}
\expandafter\gdef\csname CGBib@Grossmann1913\endcsname{J. Grossmann, "Uber die Nullstellen der Riemannschen Zeta-Funktion und der Dirichletschen L-Funktionen, doctoral thesis, Georg-August-Universit"at G"ottingen, 1913.}
\expandafter\gdef\csname CGBib@HardyLittlewoodPolya1934\endcsname{G.H. Hardy, J.E. Littlewood, G. P'olya, Inequalities, first ed., Cambridge University Press, Cambridge, 1934.}
\expandafter\gdef\csname CGBib@HasanalizadeShenWong2021\endcsname{E. Hasanalizade, Q. Shen, P.-J. Wong, Counting zeros of Dedekind zeta functions, Math. Comp. 91 (2022) 277--293, doi:10.1090/mcom/3665.}
\expandafter\gdef\csname CGBib@HasanalizadeShenWong2022\endcsname{E. Hasanalizade, Q. Shen, P.-J. Wong, Counting zeros of the Riemann zeta function, arXiv:2107.06506v1, 2021, cited with its numbering, doi:10.48550/arXiv.2107.06506; published version: J. Number Theory 235 (2022) 219--241, doi:10.1016/j.jnt.2021.06.032.}
\expandafter\gdef\csname CGBib@HiaryLeongYang2024\endcsname{G.A. Hiary, N. Leong, A. Yang, Explicit bounds for the Riemann zeta-function on the 1-line, arXiv:2306.13289v2, 2024, cited with its numbering, doi:10.48550/arXiv.2306.13289; published version: Funct. Approx. Comment. Math. 73 (2025) 53--89, doi:10.7169/facm/240711-13-8.}
\expandafter\gdef\csname CGBib@HiaryPatelYang2024\endcsname{G.A. Hiary, D. Patel, A. Yang, An improved explicit estimate for zeta(1/2+it), J. Number Theory 256 (2024) 195--217, doi:10.1016/j.jnt.2023.09.003.}
\expandafter\gdef\csname CGBib@Ivic1985\endcsname{A. Ivi'c, The Riemann Zeta-Function, Wiley, New York, 1985; reprinted by Dover, Mineola, NY, 2003.}
\expandafter\gdef\csname CGBib@Korolev2012\endcsname{M.A. Korolev, Gram's law and the argument of the Riemann zeta function, Publ. Inst. Math. (Beograd) (N.S.) 92(106) (2012) 53--78, doi:10.2298/PIM1206053K.}
\expandafter\gdef\csname CGBib@LeClair2016\endcsname{A. LeClair, Riemann hypothesis and random walks: the zeta case, arXiv:1601.00914v1, 2016, cited with its numbering; published version: Symmetry 13 (2021) 2014, doi:10.3390/sym13112014.}
\expandafter\gdef\csname CGBib@Leong2024\endcsname{N. Leong, Explicit estimates for the logarithmic derivative and the reciprocal of the Riemann zeta function, arXiv:2405.04869v5, 2026, cited with its numbering, doi:10.48550/arXiv.2405.04869.}
\expandafter\gdef\csname CGBib@LiebLoss2001\endcsname{E.H. Lieb, M. Loss, Analysis, second ed., Grad. Stud. Math. 14, Amer. Math. Soc., Providence, RI, 2001, doi:10.1090/gsm/014.}
\expandafter\gdef\csname CGBib@Mertens1874\endcsname{F. Mertens, Ein Beitrag zur analytischen Zahlentheorie, J. Reine Angew. Math. 78 (1874) 46--62, doi:10.1515/crll.1874.78.46.}
\expandafter\gdef\csname CGBib@MontgomeryVaughan1974\endcsname{H.L. Montgomery, R.C. Vaughan, Hilbert's inequality, J. Lond. Math. Soc. (2) 8 (1974) 73--82, doi:10.1112/jlms/s2-8.1.73.}
\expandafter\gdef\csname CGBib@Patel2022\endcsname{D. Patel, An explicit upper bound for |zeta(1+it)|, arXiv:2009.00769v1, 2020, cited with its numbering, doi:10.48550/arXiv.2009.00769; published version: Indag. Math. 33 (2022) 1012--1032, doi:10.1016/j.indag.2022.04.007.}
\expandafter\gdef\csname CGBib@PatelYang2023\endcsname{D. Patel, A. Yang, An explicit sub-Weyl bound for zeta(1/2+it), J. Number Theory 262 (2024) 301--334, doi:10.1016/j.jnt.2024.03.011.}
\expandafter\gdef\csname CGBib@Platt2017\endcsname{D.J. Platt, Isolating some non-trivial zeros of zeta, Math. Comp. 86 (2017) 2449--2467, doi:10.1090/mcom/3198.}
\expandafter\gdef\csname CGBib@PlattTrudgian2015\endcsname{D.J. Platt, T.S. Trudgian, An improved explicit bound on |zeta(1/2 +it)|, arXiv:1402.3953, version dated 4 July 2018, cited with its numbering, doi:10.48550/arXiv.1402.3953; published version: J. Number Theory 147 (2015) 842--851, doi:10.1016/j.jnt.2014.08.019.}
\expandafter\gdef\csname CGBib@Rademacher1959\endcsname{H. Rademacher, On the Phragm'en--Lindel"of theorem and some applications, Math. Z. 72 (1959) 192--204, doi:10.1007/BF01162949.}
\expandafter\gdef\csname CGBib@Revers2026\endcsname{M. Revers, A new improved explicit estimate for zeta(1/2+it), preprint, arXiv:2602.05614v1, 2026, doi:10.48550/arXiv.2602.05614.}
\expandafter\gdef\csname CGBib@Riesz1932\endcsname{F. Riesz, Sur un th'eor`eme de maximum de MM. Hardy et Littlewood, J. London Math. Soc. 7 (1932) 10--13, doi:10.1112/jlms/s1-7.1.10.}
\expandafter\gdef\csname CGBib@Rosser1941\endcsname{J.B. Rosser, Explicit bounds for some functions of prime numbers, Amer. J. Math. 63 (1941) 211--232, doi:10.2307/2371291.}
\expandafter\gdef\csname CGBib@Selberg1946a\endcsname{A. Selberg, Contributions to the theory of the Riemann zeta-function, Arch. Math. Naturvid. 48 (5) (1946) 89--155.}
\expandafter\gdef\csname CGBib@Tenenbaum2015\endcsname{G. Tenenbaum, Introduction to Analytic and Probabilistic Number Theory, third ed., Grad. Stud. Math. 163, Amer. Math. Soc., Providence, RI, 2015, doi:10.1090/gsm/163.}
\expandafter\gdef\csname CGBib@Titchmarsh1986\endcsname{E.C. Titchmarsh, The Theory of the Riemann Zeta-Function, second ed., revised by D.R. Heath-Brown, Clarendon Press, Oxford, 1986.}
\expandafter\gdef\csname CGBib@Trudgian2012\endcsname{T.S. Trudgian, An improved upper bound for the argument of the Riemann zeta-function on the critical line, Math. Comp. 81 (2012) 1053--1061, doi:10.1090/S0025-5718-2011-02537-8.}
\expandafter\gdef\csname CGBib@Trudgian2014\endcsname{T.S. Trudgian, An improved upper bound for the argument of the Riemann zeta-function on the critical line II, arXiv:1208.5846v2, 2012, cited with its numbering and constants, doi:10.48550/arXiv.1208.5846; published version: J. Number Theory 134 (2014) 280--292, doi:10.1016/j.jnt.2013.07.017.}
\expandafter\gdef\csname CGBib@VonMangoldt1905\endcsname{H. von Mangoldt, Zur Verteilung der Nullstellen der Riemannschen Funktion xi(t), Math. Ann. 60 (1905) 1--19, doi:10.1007/BF01447494.}
\expandafter\gdef\csname CGBib@Yang2024\endcsname{A. Yang, Explicit bounds on zeta(s) in the critical strip and a zero-free region, J. Math. Anal. Appl. 534 (2024) 128124, doi:10.1016/j.jmaa.2024.128124.}
\makeatother
